% resumo em português
\begin{resumo}
\noindent As redes ópticas passivas com arquitetura PON constituem a infraestrutura predominante no fornecimento de conectividade de alta velocidade em redes FTTH. Contudo, os processos tradicionais de dimensionamento sofrem com a fragmentação entre diagramas esquemáticos e planilhas desconectadas, resultando em falhas de cálculo no orçamento óptico e clientes inoperantes. Este trabalho tem como objetivo desenvolver uma aplicação web interativa, denominada PON Planner, que integra a modelagem visual da topologia ao cálculo automático de atenuação e sinal óptico em tempo real. A metodologia envolve a construção de um motor de cálculo em Python sob API FastAPI, integrado a uma interface reativa em React e D3.js com persistência em arquivos JSON. A solução foi validada em dois cenários reais: uma rede urbana equilibrada e uma rede rural desbalanceada composta por 14 caixas de terminação. Como resultados, o sistema respondeu ao recálculo da rede em menos de 10 ms para topologias com até 128 ONUs e identificou atenuações ocultas entre 0,7 dB e 1,3 dB omitidas no cálculo manual, reduzindo o tempo de balanceamento da rede rural de 1,5 hora para apenas 5 minutos. A principal contribuição consiste em uma ferramenta em código aberto e uma metodologia unificada que eliminam erros de dimensionamento e otimizam a viabilidade técnica de projetos FTTH.

\vspace{0.5cm}   
\noindent \textbf{Palavras-chave:} Redes Ópticas Passivas. Orçamento Óptico. PON Planner. Redes FTTH. Modelagem de Topologias.
\end{resumo}
 
% resumo em inglês
\begin{resumo}[Abstract]	
\begin{otherlanguage*}{english}
\noindent Passive optical networks with PON architecture constitute the predominant infrastructure for providing high-speed connectivity in FTTH networks. However, traditional sizing processes suffer from fragmentation between schematic diagrams and disconnected spreadsheets, resulting in optical budget calculation errors and inoperative clients. This work aims to develop an interactive web application, named PON Planner, that integrates visual topology modeling with automatic attenuation and real-time optical signal calculation. The methodology involves building a Python calculation engine under a FastAPI API, integrated into a reactive React and D3.js interface with JSON file persistence. The solution was validated across two real scenarios: a balanced urban network and an unbalanced rural network comprising 14 termination boxes. As results, the system recalculated the network in under 10 ms for topologies with up to 128 ONUs and identified hidden attenuations between 0.7 dB and 1.3 dB omitted in manual calculations, reducing the rural network balancing time from 1.5 hours to only 5 minutes. The main contribution consists of an open-source tool and a unified methodology that eliminate sizing errors and optimize the technical feasibility of FTTH projects.

\vspace{0.5cm}
\noindent \textbf{Key-words:} Passive Optical Networks. Optical Budget. PON Planner. FTTH Networks. Topology Modeling.
\end{otherlanguage*}
\end{resumo}
